% !TEX root = ../algebra_02.tex

\section{Диагонализируемые операторы. Критерий диагонализируемости в терминах геометрических кратностей}

\begin{Definition}{Диагонализируемый оператор}{}
	Оператор $\alpha\in\End V$ называется диагонализируемым, если существует базис $E$ пространства $V$, в котором матрица $[\alpha]_E$ диагональна.
\end{Definition}

\begin{Definition}{Геометрическая кратность}{}
	Если $\lambda$ --- собственное значение оператора $\alpha$, то размерность собственного подпространства $g_\lambda:=\dim V_\lambda$ называется геометрической кратностью собственного значения $\lambda$.
\end{Definition}

\begin{Lemma}{Собственные значения диагонального оператора}{}
	Пусть $E=(e_1,\ldots,e_n)$ --- базис $V$ и
	\[
	[\alpha]_E=
	\operatorname{diag}(
	\underbrace{\lambda_1,\ldots,\lambda_1}_{m_1},
	\underbrace{\lambda_2,\ldots,\lambda_2}_{m_2},
	\ldots,
	\underbrace{\lambda_k,\ldots,\lambda_k}_{m_k}
	),
	\]
	где $\lambda_1,\ldots,\lambda_k$ попарно различны.
	
	Тогда $\lambda_1,\ldots,\lambda_k$ --- все собственные значения $\alpha$, причём
	\[
	g_{\lambda_i}=m_i
	\qquad
	(i=1,\ldots,k).
	\]
\end{Lemma}

\begin{proof}
	Для любого $\lambda\in K$
	\[
	[\alpha-\lambda\operatorname{id}_V]_E
	=
	[\alpha]_E-\lambda I_n.
	\]
	Если $\lambda\notin\{\lambda_1,\ldots,\lambda_k\}$, то все диагональные элементы этой матрицы ненулевые, значит,
	$\Ker(\alpha-\lambda\operatorname{id}_V)=0$.
	Следовательно, такое $\lambda$ не является собственным значением.
	
	Пусть теперь $\lambda=\lambda_i$.
	Тогда в матрице $[\alpha-\lambda_i\operatorname{id}_V]_E$ нули стоят ровно на тех местах диагонали, которые соответствуют блоку $\lambda_i$.
	Поэтому
	\[
	V_{\lambda_i}
	=
	\Lin(e_{m_1+\cdots+m_{i-1}+1},\ldots,e_{m_1+\cdots+m_i}),
	\]
	и значит, $\dim V_{\lambda_i}=m_i$.
\end{proof}

\begin{Theorem}{Критерий через геометрические кратности}{}
	Пусть $V$ имеет размерность $n$. Оператор $\alpha$ диагонализируем тогда и только тогда, когда 
	\[
	g_{\lambda_1}+\cdots+g_{\lambda_k}=n
	\]
\end{Theorem}

\begin{proof}
	Пусть $\alpha$ диагонализируем. Тогда существует базис $E$ из собственных векторов. Разобьём этот базис на группы по собственным значениям. Если собственному значению $\lambda_i$ соответствует $m_i$ векторов базиса, то $g_{\lambda_i}=m_i$. Суммируя по всем группам, получаем общее число векторов базиса:
	\[
	g_{\lambda_1}+\cdots+g_{\lambda_k} = m_1+\cdots+m_k = n
	\].

	Обратно, пусть $g_{\lambda_1}+\cdots+g_{\lambda_k}=n$. Собственные подпространства, отвечающие разным собственным значениям, всегда образуют прямую сумму: $V_{\lambda_1}+\cdots+V_{\lambda_k} = V_{\lambda_1}\oplus\cdots\oplus V_{\lambda_k}$. Так как сумма их размерностей равна $n$, эта прямая сумма совпадает со всем пространством $V$. Выбрав базис в каждом $V_{\lambda_i}$, мы получим объединение этих базисов, которое даст базис всего $V$, состоящий исключительно из собственных векторов. Значит, $\alpha$ диагонализируем.
\end{proof}
